\subsection{Riesz spaces and Riesz systems}

\begin{definition}
  Let $X$ be a set and let $L\subset\mathbb{R}^X$ be an
  $\mathbb{R}$-vector space\footnote{
    We allow $X=\emptyset$. In that case $\mathbb{R}^X=\{\emptyset\}$ is the
    zero-dimensional vector space, and $(\emptyset,\{0\},0)$ is the only
    integration system on the empty set.
    The definitions and assertions of this paper remain meaningful in this case,
    and they can be checked directly.
  }.
  We say that $L$ is a \textbf{Riesz space} on $X$\footnote{
    What this paper calls a Riesz space is called a vector lattice in some books.
    We use ``Riesz space'' throughout.
  } if
  \begin{align*}
    f\land g &\in L \\
    f\lor g &\in L
  \end{align*}
  holds for all $f,g\in L$.
\end{definition}

\begin{remark}[on terminology]\label{rem:riesz-terminology}
  The definition above is for subspaces of $\mathbb{R}^X$, and it requires
  closure under the lattice operations with respect to the pointwise order.
  Such an $L$ is also a Riesz space in the abstract sense (an ordered vector
  space that is a lattice), with the pointwise order as its order.
  That is, a ``Riesz space on $X$'' in this paper is a special case of an
  abstract Riesz space.

  In this paper, \textbf{whenever we attach ``on $X$'' we always mean a
  sublattice of $\mathbb{R}^X$}.
  Riesz spaces in the abstract sense appear only in
  Appendix~\ref{sec:forget}, where we call them \textbf{abstract Riesz spaces}
  to keep them apart.

  This distinction is not a formality, because not every abstract Riesz space can
  be realized as a Riesz space on some $X$.
  Indeed, if $L$ is a sublattice of $\mathbb{R}^X$ then evaluation
  $\mathrm{ev}_x$ at a point gives a Riesz homomorphism from $L$ to $\mathbb{R}$,
  and these separate the elements of $L$.
  But by Proposition~\ref{prop:spec-collapses} there is no nonzero Riesz
  homomorphism $L\to\mathbb{R}$ on $L^1[0,1]$.
  Hence $L^1[0,1]$ cannot be realized as a Riesz space on any set $X$.
\end{remark}

\begin{definition}
  Let $X$ be a set and let $L$ be a Riesz space on $X$.
  The pair
  \[
    (X , L)
  \]
  is called a \textbf{Riesz system} on the set $X$.
  The set $X$ in a Riesz system $(X,L)$ is called the \textbf{base space} of the
  Riesz system $(X,L)$.
\end{definition}

\begin{definition}
  Let $(X,L)$ and $(Y,M)$ be Riesz systems on sets $X$ and $Y$
  respectively.
  A \textbf{morphism of Riesz systems}
  \[
    F:(X,L)\to(Y,M)
  \]
  is a map $\underline{F}:Y\to X$ such that
  \[
    F^*f:=f\circ\underline{F}\in M
  \]
  holds for every $f\in L$.\footnote{Strictly, a morphism is the triple
  consisting of the domain, the codomain and the map of base spaces. Below we fix
  the domain and the codomain and present a morphism by the map of base spaces.}
\end{definition}

\begin{lemma}
  \begin{enumerate}
    \item For every Riesz system $(X,L)$ on a set $X$, the identity map
      $\mathrm{id}_X=\underline{\mathrm{id}_X}:X\to X$ induces a morphism of
      Riesz systems $\mathrm{id}_{(X,L)}:(X,L)\to (X,L)$.
    \item For any two morphisms of Riesz systems $F:(X,L)\to(Y,M)$ and
      $G:(Y,M)\to(Z,N)$, the map $\underline{F}\circ\underline{G}:Z\to X$ induces
      a morphism of Riesz systems $G\circ F:(X,L)\to(Z,N)$.
  \end{enumerate}
\end{lemma}
\begin{proof}
  \begin{enumerate}
    \item Immediate, since $f\circ \operatorname{id}_X=f\in L$ for every $f\in L$.
    \item For every $f\in L$ we have $F^*f\in M$ by hypothesis. Hence, again by
      hypothesis,
      $(G\circ F)^*(f)=f\circ(\underline{F}\circ\underline{G})=(f\circ \underline{F})\circ\underline{G}\in N$,
      which is the assertion.
  \end{enumerate}
\end{proof}

\begin{definition}
  The category whose objects are Riesz systems and whose morphisms are morphisms
  of Riesz systems is written
  \[
    \mathbf{RieszSys}
  \]
  and called the \textbf{category of Riesz systems}.
\end{definition}
